\begin{summarybox}
  \textbf{\textcolor{CSummary}{一句话记住}}
  \textbf{上下平移，改的是高度：上移整体加，下移整体减；$x$保持不变。}
  \[
    \boxed{\text{上移 }b:\ y=f(x)+b}\qquad
    \boxed{\text{下移 }b:\ y=f(x)-b}\quad(b>0).
  \]

  \medskip
  \textbf{\textcolor{CSummary}{考场三步（快用）}}
  \begin{enumerate}[leftmargin=*,itemsep=0.22em,topsep=0.2em]
    \item \textbf{看方向：}上移写“$+$”，下移写“$-$”。
    \item \textbf{看位置：}把数加减在\textbf{整个$f(x)$外面}；不要随意放进$x$的括号里。
    \item \textbf{看结果：}同一个$x$，新旧函数值恰好相差平移量；定义域不变。
  \end{enumerate}

  \textbf{\textcolor{CSummary}{最短例子}}
  $y=x^2\xrightarrow{\text{上移 }2}y=x^2+2$；
  $y=x^2\xrightarrow{\text{下移 }3}y=x^2-3$。

  \medskip
  \textbf{\textcolor{CSummary}{特别注意}}
  $y=(x+2)^2$是向\textbf{左}移$2$，不是向上移$2$。连续上下移动可把变化量相加，例如先上移$1$再下移$3$，相当于\textbf{下移$2$}。
\end{summarybox}
